\documentclass[a4paper,12pt]{article}
\usepackage[utf8]{inputenc}
\usepackage[T1]{fontenc}
\usepackage[ngerman]{babel}
\usepackage{geometry}
\geometry{margin=2.5cm}
\begin{document}

\section*{Imaging and Seismic Modelling Inside Volcanoes Using Machine Learning}
\textbf{Autoren:} Gareth Shane O'Brien, Christopher J. Bean, Hugo Meiland, Philipp Witte \\
\textbf{Journal:} Scientific Reports, 13:630 \\
\textbf{Jahr:} 2023

\subsection*{Zusammenfassung}

Die Arbeit adressiert die grundlegende Herausforderung der seismischen Bildgebung in vulkanischen 
Umgebungen, wo starke strukturelle Heterogenitäten seismische Wellen stark streuen und konventionelle 
tomographische Methoden scheitern. Passive Erdbebentomographie liefert lediglich niederfrequente, 
geglättete Geschwindigkeitsmodelle, die die feinskaligen geologischen Strukturen nicht auflösen können, 
welche für die starken Streupfadeffekte verantwortlich sind. Als neuartiger Ansatz wird ein 
\textit{Fourier Neural Operator} (FNO) eingesetzt, der sowohl das Vorwärtsproblem der 
Wellenfeldsimulation als auch das inverse Problem der Geschwindigkeitsinversion löst.

\subsubsection*{Datengenerierung und Trainingsaufbau}

Da keine physikalischen Datensätze mit bekanntem Untergrundmodell für das Training geeignet sind, 
werden synthetische Trainingsdaten durch numerische Simulation erzeugt. Als Simulator wird eine 
2D-Gittermethode für elastische Wellenausbreitung verwendet, die breitbandige Heterogenitäten und 
starke Geschwindigkeitsgradienten korrekt behandelt. Zwei verschiedene Populationen von 
Geschwindigkeitsmodellen werden definiert: Population~L mit geschichteten Modellen und 
Population~T mit tomographieartig geglätteten Perturbationen. Jede Population umfasst 20.000 
individuelle Modelle mit zufälligen Quellpositionen und Quellmechanismen, was zu insgesamt 
über 40.000 Simulationen für das FNO-Training führt. Eine unabhängige Validierungspopulation~V 
mit stärkeren Perturbationen dient zur Überprüfung der Generalisierungsfähigkeit.

\subsubsection*{FNO-Architektur und Training}

Der FNO ist eine neuartige Netzwerkarchitektur, die partielle Differentialgleichungen als Operatoren 
auffasst und lernt, diese Operatoren -- sowie ihre Inversen -- auf Funktionenräumen abzubilden. 
Im Gegensatz zu klassischen Faltungsnetzwerken (CNN) arbeitet der FNO im Fourierraum, was 
drei wesentliche Vorteile bietet: (1) globale Erfassung der Kontinuität von PDEs gegenüber 
lokalen Faltungskernen, (2) Auflösungs- und Gitterinvarianz, sodass mit niedrig aufgelösten 
Trainingsdaten höher aufgelöste Vorhersagen möglich sind, und (3) erhöhte Recheneffizienz 
durch Trunkierung hoher Frequenzen. Das Netzwerk ist vier Schichten tief, in PyTorch implementiert 
und wird mit AdamW-Optimierer über 200 Epochen mit variabler Lernrate trainiert.

Für das \textit{Vorwärtsproblem} erhält der FNO das Geschwindigkeitsmodell, die Quellposition 
und den Quellmechanismus als Eingabe und gibt das elastische Wellenfeld bzw. synthetische 
Seismogramme aus. Für das \textit{inverse Problem} werden Eingaben und Ausgaben vertauscht: 
zehn zufällig ausgewählte zweikomponentige Seismogramme (x- und z-Komponente) dienen als 
Eingabe, während das vorhergesagte Geschwindigkeitsmodell und die Quellposition als Ausgabe 
geliefert werden. Dem Training werden 15\,\% weißes Rauschen auf den Eingabespuren hinzugefügt, 
um die Robustheit gegenüber realen Messbedingungen zu erhöhen.

\subsubsection*{Ergebnisse}

Das trainierte Vorwärtsmodell reproduziert die dominanten Phasen des elastischen Wellenfelds -- 
inklusive P- und S-Wellen -- mit hoher Genauigkeit, wenngleich kleinamplitudigere 
Streuereignisse teilweise nicht vollständig erfasst werden. Die Inversionsresultate zeigen, 
dass der FNO aus lediglich einem Schussexperiment (10 Oberflächen-Empfänger) ein 
hochaufgelöstes Geschwindigkeitsmodell rekonstruieren kann, das weit über die 
Auflösungsfähigkeiten konventioneller Tomographie hinausgeht. Quantifiziert anhand des 
normierten RMS-Fehlers über 20 ungesehene Testmodelle aus beiden Populationen zeigen die 
Geschwindigkeitsmodelle eine hohe Vorhersagegenauigkeit, während die Quelllokalisation 
empfindlicher gegenüber Rauschen reagiert. Die Generalisierung auf die abweichende 
Population~V gelingt durch Stapelung von Vorhersagen aus 100 Schussexperimenten, was 
die Konvergenz zum wahren Modell klar demonstriert.

\subsubsection*{Diskussion und Ausblick}

Die Autoren betonen, dass der FNO keine traditionellen PDE-Löser ersetzen soll -- diese 
werden für die Trainingsdatengenerierung benötigt -- sondern als schnelles Surrogat-Modell 
für Vorwärts- und Inversionsaufgaben dient. Die Skalierung auf 3D wird als technisch 
machbar, jedoch rechenintensiv eingestuft, wobei die Auflösungsinvarianz des FNO durch 
Verwendung niedrig aufgelöster Trainingsdaten helfen kann. Die zentrale offene Frage 
bleibt, wie eine Trainingsmodellpopulation konstruiert werden muss, die hinreichend 
repräsentativ für die reale vulkanische Struktur ist, um eine Inversion mit echten 
Felddaten zu ermöglichen.

\subsection*{Bezug zur Masterarbeit}

Diese Arbeit weist mehrere direkte Parallelen zur vorliegenden Masterarbeit über 
physik-basiertes maschinelles Lernen zur Vorhersage von Strömungsfeldern um sinkende 
Kristalle auf.

\paragraph{Gemeinsame Architektur.}
Beide Arbeiten setzen den FNO als zentrale Netzwerkarchitektur ein. Während O'Brien 
et al.\ den FNO zur Lösung der elastischen Wellengleichung nutzen, wird er in der 
Masterarbeit als Surrogat-Modell für die Stokes-Gleichungen in magmatischen 
Sedimentationssystemen eingesetzt. Die strukturellen Vorteile des FNO -- globale 
Spektralverarbeitung, Auflösungsinvarianz und Recheneffizienz -- sind in beiden 
Domänen gleichermaßen relevant.

\paragraph{Surrogat-Modellierung physikalischer Simulationen.}
In beiden Fällen dienen numerische PDE-Löser als Ground-Truth-Datengeneratoren 
(elastische Gittermethode bzw. LaMEM), und das neuronale Netz lernt, diese Löser 
mit drastisch reduziertem Rechenaufwand zu approximieren. Die Masterarbeit folgt 
damit demselben Paradigma: Das trainierte Netz ist nicht dazu gedacht, den Simulator 
zu ersetzen, sondern ermöglicht schnelle Vorhersagen für Konfigurationen, die ohne 
Surrogat-Modell prohibitiv teuer wären.

\paragraph{Generalisierung als zentrale Herausforderung.}
O'Brien et al.\ identifizieren die Generalisierung auf Modellpopulationen außerhalb 
der Trainingsverteilung als Schlüsselproblem und lösen es partiell durch Stapelung 
von Vorhersagen. Dies entspricht direkt der zentralen Forschungsfrage der Masterarbeit: 
Kann ein auf $N$-Kristall-Konfigurationen trainiertes Netzwerk auf Konfigurationen 
mit $N{+}m$ Kristallen generalisieren? Beide Arbeiten thematisieren damit die 
fundamentale Extrapolationsschwäche neuronaler Netze und entwickeln Strategien, 
um diese zu überwinden.

\paragraph{Physikalisch motivierte Trainingsdatenstrategie.}
Beide Arbeiten verwenden statistisch parametrisierte Modellpopulationen, um den 
Trainingsraum zu definieren. O'Brien et al.\ schlagen vor, geologisches Vorwissen 
(Topographie, Petrophysik, Bohrlochsdaten) zur Einschränkung der Trainingsverteilung 
zu nutzen -- ein Prinzip, das in der Masterarbeit durch physikalisch motivierte 
Normierungsstrategien und Curriculum-Learning (von Einzel- zu Mehrkristall-Konfigurationen) 
umgesetzt wird.

\paragraph{Abgrenzung.}
Im Unterschied zur vorliegenden Masterarbeit behandelt O'Brien et al.\ ein 
Wellenpropagationsproblem mit zeitabhängigen, transienten Lösungen. Die Masterarbeit 
hingegen fokussiert auf stationäre Stokes-Strömungen, was die Verwendung von 
Stromfunktionsansätzen ermöglicht, die die Divergenzfreiheit der Strömung 
mathematisch garantieren -- ein Aspekt, der in der seismischen FNO-Anwendung 
keine Entsprechung findet. Zudem wird in der Masterarbeit explizit der Residuallernansatz 
$\mathbf{v}_\mathrm{total} = \mathbf{v}_\mathrm{Stokes,analytisch} + \Delta\mathbf{v}_\mathrm{gelernt}$ 
verfolgt, der durch Ausnutzung analytischer Einzelkristalllösungen die zu lernende 
Funktion vereinfacht.

\end{document}